% !TeX root = ../../Book2.tex
\section{多线性 Fourier 限制}
\label{b2:sec:2504}

完整 \(d\) 线性横截问题是 Loomis–Whitney 模型的振荡推广。平直情形已经证明；一般曲面允许法向变化，需要尺度归纳和多线性管束理论。

\subsection{多线性延拓估计}
\label{b2:subsec:250401}

下面采用\term[bennett-carbery-tao-duoxianxing-xianzhi]{Bennett--Carbery--Tao 多线性限制定理}[Bennett--Carbery--Tao multilinear restriction theorem]的下述版本。

\begin{theorem}[深层结构定理：Bennett--Carbery--Tao 局部多线性限制]\label{b2:thm:25-bct}
设 \(d\geq2\)，\(\Sigma_j:U_j\subset\mathbb R^{d-1}\to\mathbb R^d\) 为固定紧参数片上的光滑参数化，\(C^2\) 范数有统一上界，切向楔积 \(Y_j\) 满足
\(|\det(Y_1(u_1),\ldots,Y_d(u_d))|\ge\nu>0\)。
令 \(E_jg_j(x)=\int_{U_j}e^{2\pi\mathrm ix\cdot\Sigma_j(u)}g_j(u)\,\mathrm du\)。
对每个 \(\epsilon>0\)、\(R\ge1\)，有
\[
 \left\|\prod_{j=1}^dE_jg_j\right\|_{L^{2/(d-1)}(B_R)}
       \le C_\epsilon R^\epsilon\prod_j\|g_j\|_2.
\]
\end{theorem}
\begin{proofsketch}
波包分解将延拓转化为横截管束的组合；尺度归纳把限制界归约到多线性 Kakeya 控制，再以 Gaussian 族的单调性处理管束。误差和尺度迭代产生允许的 \(R^\epsilon\) 损失。
完整管束与归纳证明在此略去。上述版本是
\cite[Theorem 1.16, p.~267; Sections 2--4]{Bennett2006Multilinear}
的 \(L^2\) 端点特例；正式期刊版的假设、指数及尺度损失已核验。
\end{proofsketch}
保留局部球和 \(R^\epsilon\) 很重要。本章没有证明去掉这个因子的全空间端点界。平直模型能无损成立，不等于一般曲面也可直接无损；同样不能把小于1的输出准范数当作 Banach 范数随意使用三角不等式。

\subsection{多线性限制}
\label{b2:subsec:250402}

多线性限制通常用延拓乘积陈述。它与单个线性算子的简单伴随关系不同；固定其他输入后可以得到一个线性配对，再求其伴随，但完整形式涉及多个曲面。
参数测度 \(\mathrm du\) 与面测度的关系为
\(\mathrm d\sigma_j=|Y_j(u)|\,\mathrm du\)。
若参数化的 Jacobi 因子有正上下界，转换输入范数只改变常数；若没有这个下界，不能省略该因子。

\ProseParagraphBreak
最小维数 \(d=2\) 时乘积输出为 \(L^2\)，上一节的曲线和映射证明已给出局部横截模型。高维输出指数 \(2/(d-1)\) 越来越小，反映的是多个独立方向同时使乘积集中区域缩小，而非单个因子的可积性无限改善。

\subsection{对线性限制的反馈}
\label{b2:subsec:250403}

线性输出在一个点很大时，可能主要来自一个帽，或来自多个法向分离的帽。前一种情况可在更小尺度重新规范化，后一种情况可以调用多线性估计。这种“分离方向部分加集中部分”的组织让横截界反馈到线性限制。
实际归约需证明每个点的帽分配、范数求和与尺度损失的递推；仅知道多线性定理不自动给出线性猜想全部范围。

\ProseParagraphBreak
横截 determinant 接近零的帽族不能硬塞进 \(d\) 线性定理。它们可能集中在较低维方向簇中，需要双线性、低维或切向方法。分辨哪些输入属于横截部分，是现代归纳与分割的关键步骤。
